% =========================================================================
% secciones_zh/sec06_discussion_policy.tex — 第VI节，中文版
% =========================================================================
\section{讨论与政策启示}
\label{sec:discussion}

\textbf{面向飓风暴露岛屿的确定性韧性。}加勒比岛屿系统的规划对象不是统计意义上的天气，而是有名字的风暴。一个其 $N-1$ 裕度携带相关性诚实自助包络、其最大可信预想事故被显式演练、其状态分辨率从容处于半周期之内的基准，把韧性从愿景转变为可排程的计算。对一个十二个月内两次因全国性停电而全面动员的国家\cite{ocseni}而言，耗时数小时的传统研究与2.65~ms国家级状态之间的差别，就是事后重建事件与事前预判事件之间的差别。

\textbf{作为计算边界的规划法。}本基准的划分不是优化产物：它就是国家领土法令\cite{ley368}。当电气划分与第368-22号法律的规划区重合时，每个仿真结果都已然以审批、土地利用与微电网自治决策所用的单位表达——自治问题（失去一条干线后哪片领土可以独立支撑？）由此成为模型的属性，而非另立的课题。可再生能源激励框架\cite{ley57}与国家扩张规划周期\cite{cne2022}由此获得共同的证据基底。

\textbf{标准对齐与数据主权。}本基准以其必须互操作的标准语言表述自身：CIM一致的模型交换\cite{iec61970}、测量级的电能质量定义\cite{iec61000}、\cite{ieee1459}、谐波实践\cite{ieee519}、研究方法推荐实践\cite{ieee399}，以及逆变器型资源互联\cite{ieee2800}、\cite{ieee1547}。与此同时，底层运行档案保留在国家托管之下：已发布基准携带结构、参数与验证工具，而主权遥测层由国家机构治理。这一分离——结构上的开放科学、档案上的主权托管——是刻意的政策架构，而非疏漏。

\textbf{能力建设。}可以被开放仿真的国家电网，就是可以被开放教授的国家电网。服务于本文的同一工件也服务于高校实验室：博士与本科练习面对真实资产执行——它的距离、它的停运、它的锚点——而非1962年的制图台式抽象，从而弥合教室与控制室之间的鸿沟。

\textbf{局限。}本次发布的诚实边界为：（i）69~kV层经静态等效进入，对潮流在边界处精确，但对凝聚区域内部的短路暂态并不精确；（ii）列为相量域列式——电磁暂态（EMT）研究需要EMT扩展；（iii）遥测以运行时延到达，孪生是因果的而非先知的；（iv）市场结算不在本次发布范围之内，尽管水平域的小时分辨率已达结算级。

\textbf{向其他岛屿系统的复制。}最后，本方法论是刻意可移植的，移植协议本身即是贡献的一部分。由于不变列式不触及任何地理，为另一岛屿系统复制SENI-142只需重新实例化领土条款：（1）按同一节点规则组装电网资产（每变电站每电压等级一个节点，变压器为支路）；（2）在同一数据协议语句下摄取运行水平域，并指名可替换的系统运营商；（3）从该法域的治理性领土法令声明划分——无论其为何——并按干线联络拓扑将法域层捆绑为电气宏域；（4）从第一道校验门起重新封存确定性生命周期。该共和国的基准由此同时充当工件与模板——既是一国自己的仪器，也是给岛屿世界的可复现性契约。
